\begin{textsample}{POS Dim 4 – 2018 – Score 18.00 – 2018/t000514.txt}  \label{ex:f4_pos_064}
It was my first year as an atmospheric science professor at Texas Tech University.
We had just moved to Lubbock, Texas, which had recently been named the second most \textbf{conservative} city in the entire United States.
A colleague asked me to guest teach his undergraduate geology class.
I said,"Sure."But when I showed up, the lecture hall was cavernous and dark.
As I tracked the history of the carbon cycle through geologic time to present day, most of the students were slumped over, dozing or looking at their \textbf{phones}.
I ended my talk with a hopeful request for any questions.
And one hand shot up right away.
I looked encouraging, he stood up, and in a loud voice, he said,"You ' re a democrat, aren ' t you?""No,"I said,"I ' m Canadian."That was my baptism by fire into what has now become a sad fact of life here in the United States and increasingly across Canada as well.
The fact that the number one predictor of whether we agree that \textbf{climate} is changing, humans are responsible and the \textbf{impacts} are increasingly serious and even dangerous, has nothing to do with how much we know about science or even how smart we are but simply where we fall on the political spectrum.
Does the thermometer give us a different answer depending on if we ' re liberal or \textbf{conservative}?
Of course not.
But if that thermometer tells us that the planet is warming, that humans are responsible and that to fix this thing, we have to wean ourselves off fossil fuels as soon as possible — well, some people would rather cut off their arm than give the government any further excuse to disrupt their comfortable lives and tell them what to do.
But saying,"Yes, it ' s a real problem, but I don ' t want to fix it,"that makes us the bad \textbf{guy}, and nobody wants to be the bad \textbf{guy}.
So instead, we use arguments like,"It ' s just a natural cycle.""It ' s the sun."Or my favorite,"Those \textbf{climate} scientists are just in it for the money."I get that at least once a week.
But these are just sciencey-sounding smoke screens, that are designed to hide the real reason for our objections, which have nothing to do with the science and everything to do with our ideology and our identity.
So when we turn on the TV these days, it seems like pundit X is saying,"It ' s cold outside.
Where is global warming now?"And politician Y is saying,"For every scientist who says this thing is real, I can find one who says it isn ' t."So it ' s no surprise that sometimes we feel like \textbf{everybody} is saying these myths.
But when we look at the data — and the Yale Program on \textbf{Climate} [ Change ] Communication has done public opinion polling across the country now for a number of years — the data shows that actually 70 percent of people in the United States agree that the \textbf{climate} is changing.
And 70 percent also agree that it will harm plants and animals, and it will harm future generations.
But then when we dig down a bit deeper, the rubber starts to hit the road.
Only about 60 percent of people think it will affect people in the United States.
Only 40 percent of people think it will affect us personally.
And then when you ask people,"Do you ever talk about this?"two-thirds of people in the entire United States say,"Never."And even worse, when you say,"Do you hear the media talk about this?"Over three-quarters of people say no.
So it ' s a vicious cycle.
The planet warms.
Heat waves get stronger.
Heavy precipitation gets more frequent.
Hurricanes get more intense.
Scientists release yet another doom-filled report.
Politicians push back even more strongly, repeating the same sciencey-sounding myths.
What can we do to break this vicious cycle?
The number one thing we can do is the exact thing that we ' re not doing : talk about it.
But you might say,"I ' m not a scientist.
How am I supposed to talk about radiative forcing or cloud parametrization in \textbf{climate} models?"We don ' t need to be talking about more science ; we ' ve been talking about the science for over 150 years.
Did you know that it ' s been 150 years or more since the 1850s, when \textbf{climate} scientists first discovered that digging up and burning coal and gas and oil is producing heat- trapping gases that is wrapping an extra blanket around the planet?
That ' s how long we ' ve known.
It ' s been 50 years since scientists first formally warned a US president of the dangers of a changing \textbf{climate}, and that president was Lyndon B.
Johnson.
And what ' s more, the social science has taught us that if people have built their identity on rejecting a certain set of facts, then arguing over those facts is a personal attack.
It causes them to dig in deeper, and it digs a trench, rather than building a bridge.
So if we aren ' t supposed to talk about more science, or if we don ' t need to talk about more science, then what should we be talking about?
The most important thing to do is, instead of starting up with your head, with all the data and facts in our head, to start from the heart, to start by talking about why it matters to us, to begin with genuinely shared values.
Are we both parents?
Do we live in the same community?
Do we enjoy the same outdoor activities : hiking, biking, fishing, even hunting?
Do we care about the economy or national security?
For me, one of the most foundational ways I found to connect with people is through my faith.
As a Christian, I believe that God created this incredible planet that we live on and gave us responsibility over every living thing on it.
And I furthermore believe that we are to care for and love the least fortunate among us, those who are already suffering the \textbf{impacts} of poverty, hunger, disease and more.
If you don ' t know what the values are that someone has, have a conversation, get to know them, figure out what makes them tick.
And then once we have, all we have to do is connect the dots between the values they already have and why they would care about a changing \textbf{climate}.
I truly believe, after thousands of conversations that I ' ve had over the past decade and more, that just about every single person in the world already has the values they need to care about a changing \textbf{climate}.
They just haven ' t connected the dots.
And that ' s what we can do through our conversation with them.
The only reason why I care about a changing \textbf{climate} is because of who I already am.
I ' m a mother, so I care about the future of my child.
I live in West Texas, where water is already scarce, and \textbf{climate} change is \textbf{impacting} the availability of that water.
I ' m a Christian, I care about a changing \textbf{climate} because it is, as the military calls it, a"threat multiplier."It takes those issues, like poverty and hunger and disease and lack of access to clean water and even political \textbf{crises} that lead to \textbf{refugee} \textbf{crises} — it takes all of these issues and it exacerbates them, it makes them worse.
I ' m not a Rotarian.
But when I gave my first talk at a Rotary Club, I walked in and they had this giant banner that had the Four-Way Test on it.
Is it the truth?
Absolutely.
Is it fair?
Heck, no, that ' s why I care most about \textbf{climate} change, because it is absolutely unfair.
Those who have contributed the least to the problem are bearing the brunt of the \textbf{impacts}.
It went on to ask : Would it be beneficial to all, would it build goodwill?
Well, to fix it certainly would.
So I took my talk, and I reorganized it into the Four-Way Test, and then I gave it to this group of \textbf{conservative} businesspeople in West Texas.
And I will never forget at the end, a local bank owner came up to me with the most bemused look on his face.
And he said,"You know, I wasn ' t sure about this whole global warming thing, but it passed the Four-Way Test."These values, though — they have to be genuine.
I was giving a talk at a Christian college a number of years ago, and after my talk, a fellow scientist came up and he said,"I need some help.
I ' ve been really trying hard to get my foot in the door with our local churches, but I can ' t seem to get any traction.
I want to talk to them about why \textbf{climate} change matters."So I said,"Well, the best thing to do is to start with the denomination that you ' re part of, because you share the most values with those people.
What type of church do you attend?""Oh, I don ' t attend any church, I ' m an atheist,"he said.
I said,"Well, in that case, starting with a faith community is probably not the best idea.
Let ' s talk about what you do enjoy doing, what you are \textbf{involved} in."And we were able to identify a community group that he was part of, that he could start with.
The bottom line is, we don ' t have to be a liberal tree hugger to care about a changing \textbf{climate}.
All we have to be is a human living on this planet.
Because no matter where we live, \textbf{climate} change is already affecting us today.
If we live along the coasts, in many places, we ' re already seeing"sunny-day flooding."If we live in western North America, we ' re seeing much greater area being burned by wildfires.
If we live in many coastal locations, from the Gulf of Mexico to the South Pacific, we are seeing stronger hurricanes, typhoons and cyclones, powered by a warming ocean.
If we live in Texas or if we live in Syria, we ' re seeing \textbf{climate} change supersize our droughts, making them more frequent and more severe.
Wherever we live, we ' re already being affected by a changing \textbf{climate}.
So you might say,"OK, that ' s good.
We can talk \textbf{impacts}.
We can scare the pants off people, because this thing is serious."And it is, believe me.
I ' m a scientist, I know.
But fear is not what is going to motivate us for the long-term, sustained change that we need to fix this thing.
Fear is designed to help us run away from the bear.
Or just run faster than the person beside us.
What we need to fix this thing is rational hope.
Yes, we absolutely do need to recognize what ' s at stake.
Of course we do.
But we need a vision of a better future — a future with abundant energy, with a stable economy, with resources available to all, where our lives are not worse but better than they are today.
There are solutions.
And that ' s why the second important thing that we have to talk about is solutions — practical, viable, accessible, attractive solutions.
Like what?
Well, there ' s no silver bullet, as they say, but there ' s plenty of silver buckshot.
There ' s simple solutions that save us money and reduce our carbon footprint at the same time.
Yes, light bulbs.
I love my plug-in car.
I ' d like some solar shingles.
But imagine if every home came with a switch beside the front door, that when you left the house, you could turn off everything except your fridge.
And maybe the DVR.
Lifestyle choices : eating local, eating lower down the food chain and reducing food waste, which at the global scale, is one of the most important things that we can do to fix this problem.
I ' m a \textbf{climate} scientist, so the irony of traveling around to talk to people about a changing \textbf{climate} is not lost on me.
The biggest part of my personal carbon footprint is my travel.
And that ' s why I carefully collect my invitations.
I usually don ' t go anywhere unless I have a critical mass of invitations in one place — anywhere from three to four to sometimes even as many as 10 or 15 talks in a given place — so I can minimize the \textbf{impact} of my carbon footprint as much as possible.
And I ' ve transitioned nearly three-quarters of the talks I give to video.
Often, people will say,"Well, we ' ve never done that before."But I say,"Well, let ' s give it a try, I think it could work."Most of all, though, we need to talk about what ' s already happening today around the world and what could happen in the future.
Now, I live in Texas, and Texas has the highest carbon emissions of any state in the United States.
You might say,"Well, what can you talk about in Texas?"The answer is : a lot.
Did you know that in Texas there ' s over 25, 000 jobs in the wind energy industry?
We are almost up to 20 percent of our electricity from clean, renewable sources, most of that wind, though solar is growing quickly.
The largest army base in the United States, Fort Hood, is, of course, in Texas.
And they ' ve been powered by wind and solar energy now, because it ' s saving taxpayers over 150 million dollars.
Yes.
What about those who don ' t have the resources that we have?
In sub-Saharan Africa, there are hundreds of millions of people who don ' t have access to any type of energy except kerosine, and it ' s very expensive.
Around the entire world, the fastest-growing type of new energy today is solar.
And they have plenty of solar.
So social \textbf{impact} investors, nonprofits, even corporations are going in and using innovative new microfinancing schemes, like, pay-as-you-go solar, so that people can buy the power they need in increments, sometimes even on their cell \textbf{phone}.
One company, Azuri, has distributed tens of thousands of units across 11 countries, from Rwanda to Uganda.
They estimate that they ' ve powered over 30 million hours of electricity and over 10 million hours of cell \textbf{phone} charging.
What about the giant growing economies of China and India?
Well, \textbf{climate} \textbf{impacts} might seem a little further down the road, but \textbf{air} quality \textbf{impacts} are right here today.
And they know that clean energy is essential to powering their future.
So China is \textbf{investing} hundreds of billions of dollars in clean energy.
They ' re flooding coal mines, and they ' re putting floating solar panels on the surface.
They also have a panda-shaped solar farm.
Yes, they ' re still burning coal.
But they ' ve shut down all the coal plants around Beijing.
And in India, they ' re looking to replace a quarter of a billion incandescent light bulbs with LEDs, which will save them seven billion dollars in energy costs.
They ' re investing in green jobs, and they ' re looking to decarbonize their entire vehicle fleet.
India may be the first country to industrialize without relying primarily on fossil fuels.
The world is changing.
But it just isn ' t changing fast enough.
Too often, we picture this problem as a giant boulder sitting at the bottom of a hill, with only a few hands on it, trying to roll it up the hill.
But in reality, that boulder is already at the top of the hill.
And it ' s got hundreds of millions of hands, maybe even billions on it, pushing it down.
It just isn ' t going fast enough.
So how do we speed up that giant boulder so we can fix \textbf{climate} change in time?
You guessed it.
The number one way is by talking about it.
The bottom line is this : \textbf{climate} change is affecting you and me right here, right now, in the places where we live.
But by working together, we can fix it.
Sure, it ' s a daunting problem.
Nobody knows that more than us \textbf{climate} scientists.
But we can ' t give in to despair.
We have to go out and actively look for the hope that we need, that will inspire us to act.
And that hope begins with a conversation today.
Thank you.

% matched lemmas: air, climate, conservative, crisis, everybody, guy, impact, invest, involve, phone, refugee
\end{textsample}
